% (c) Nikita Lisitsa, lisyarus@gmail.com, 2026

\documentclass[10pt]{beamer}

\usepackage[T2A]{fontenc}
\usepackage[russian]{babel}
\usepackage{minted}

\usepackage[most]{tcolorbox}
\tcbuselibrary{minted}

\usepackage{graphicx}
\graphicspath{ {./images/} }

\usepackage{adjustbox}

\usepackage{color}
\usepackage{soul}
\usepackage{multicol}
\usepackage{hyperref}
\usepackage{tikz}
\usetikzlibrary{arrows.meta,positioning,calc}

\usetheme{metropolis}

\definecolor{red}{rgb}{1,0,0}
\definecolor{codebg}{RGB}{29,35,49}
\setminted{fontsize=\footnotesize}

\makeatletter
\newcommand{\slideimage}[1]{
  \begin{figure}
    \begin{adjustbox}{width=\textwidth, totalheight=\textheight-2\baselineskip-2\baselineskip,keepaspectratio}
      \includegraphics{#1}
    \end{adjustbox}
  \end{figure}
}
\makeatother

\title{Язык программирования C++}
\subtitle{Лекция 18: Приведение типов, RTTI, typeid, lvalue/rvalue, move, вывод типов, range-based for}
\date{2026}
\author{Никита Лисица}

\setbeamertemplate{footline}[frame number]

\begin{document}

\frame{\titlepage}

\begin{frame}[fragile]
\frametitle{Напоминание: приведение типов в C++}
\begin{itemize}
\usemintedstyle{tango}
\item Указатели автоматически приводятся к \mintinline{cpp}|void*|, к константной версии того же указателя, и к указателю на родительский класс:
\usemintedstyle{lightbulb}
\begin{tcblisting}{listing engine=minted, minted language=cpp, minted options={fontsize=\scriptsize}, listing only, colback=codebg, colframe=codebg, rounded corners}
int i = 12;
void* p = &i; // OK
const int* q = &i; // OK

struct A {};
struct B : A {};

B b;
A* pa = &b; // OK
\end{tcblisting}
\end{itemize}
\end{frame}

\begin{frame}[fragile]
\frametitle{Напоминание: приведение типов в C++}
\begin{itemize}
\usemintedstyle{tango}
\item В других случаях будет ошибка компиляции:
\usemintedstyle{lightbulb}
\begin{tcblisting}{listing engine=minted, minted language=cpp, minted options={fontsize=\scriptsize}, listing only, colback=codebg, colframe=codebg, rounded corners}
void* p = ...;
int* q1 = p; // Ошибка компиляции
int* q2 = (int*)p; // OK, явное приведение

const int i = 12;
int* pi1 = &i; // Ошибка компиляции
int* pi2 = (int*)&i; // OK, явное приведение

A a;
B* pb1 = &a; // Ошибка компиляции
B* pb2 = (B*)&a; // OK, явное приведение
\end{tcblisting}
\end{itemize}
\end{frame}

\begin{frame}[fragile]
\frametitle{Напоминание: приведение типов в C++}
\begin{itemize}
\usemintedstyle{tango}
\item Все числовые типы приводятся друг к другу:
\usemintedstyle{lightbulb}
\begin{tcblisting}{listing engine=minted, minted language=cpp, minted options={fontsize=\scriptsize}, listing only, colback=codebg, colframe=codebg, rounded corners}
int i = 12;
float f = i;
double d = f;
char c = d;
std::uint32_t u = f;
\end{tcblisting}
\pause
\item \textbf{NB}: часть таких приведений -- \textit{undefined behavior}
\end{itemize}
\end{frame}

\begin{frame}[fragile]
\frametitle{Напоминание: приведение типов в C++}
\begin{itemize}
\usemintedstyle{tango}
\item Для приведения своих типов (структур/классов) можно сделать неявный конструктор или оператор приведения:
\usemintedstyle{lightbulb}
\begin{tcblisting}{listing engine=minted, minted language=cpp, minted options={fontsize=\scriptsize}, listing only, colback=codebg, colframe=codebg, rounded corners}
class BigInteger {
public:
  BigInteger(int value); // Приведение int -> BigInteger
  operator int() const; // Приведение BigInteger -> int
};
\end{tcblisting}
\pause
\usemintedstyle{tango}
\item \textbf{NB}: если для приведения типов \mintinline{cpp}|A -> B| есть как оператор приведения \mintinline{cpp}|A::operator B()|, так и конструктор \mintinline{cpp}|B::B(A)|, будет ошибка компиляции
\end{itemize}
\end{frame}

\begin{frame}[fragile]
\frametitle{Приведение типов в C++}
\begin{itemize}
\usemintedstyle{tango}
\item Конструкция \mintinline{cpp}|(int*)p| или \mintinline{cpp}|(float)42| или \mintinline{cpp}|(int)myBigInt| называется \textit{приведением типов в стиле Си (C-style cast)}
\pause
\item В C++ принято её не использовать, так как она умеет слишком многое, и плохо выражает намерение программиста
\pause
\item Вместо этого, есть 4 специализированных оператора приведения типов, которые выглядят как шаблонные функции:
\pause
\begin{itemize}
\item \mintinline{cpp}|static_cast<T>(x)|
\pause
\item \mintinline{cpp}|reinterpret_cast<T>(x)|
\pause
\item \mintinline{cpp}|const_cast<T>(x)|
\pause
\item \mintinline{cpp}|dynamic_cast<T>(x)|
\end{itemize}
\end{itemize}
\end{frame}

\begin{frame}[fragile]
\frametitle{static\_cast}
\begin{itemize}
\usemintedstyle{tango}
\item Всё, что умеют неявные приведения типов
\pause
\item Приведение указателя на базовый класс к указателю на наследник \textit{(downcasting)}
\pause
\item Приведение указателя на \mintinline{cpp}|void| к другому типу указателя
\pause
\usemintedstyle{lightbulb}
\begin{tcblisting}{listing engine=minted, minted language=cpp, minted options={fontsize=\scriptsize}, listing only, colback=codebg, colframe=codebg, rounded corners}
int i = static_cast<int>(3.14f);

A* pa = new B;
B* pb = static_cast<B*>(pa);

void* pv = pa;
A* pa2 = static_cast<A*>(pv);
\end{tcblisting}
\pause
\item Легче увидеть в коде, понятнее намерения программиста
\end{itemize}
\end{frame}

\begin{frame}[fragile]
\frametitle{reinterpret\_cast}
\begin{itemize}
\usemintedstyle{tango}
\item Приведение между указателями любых типов
\pause
\item Приведение между ссылками любых типов
\pause
\item Приведение между указателями и числами
\pause
\usemintedstyle{lightbulb}
\begin{tcblisting}{listing engine=minted, minted language=cpp, minted options={fontsize=\scriptsize}, listing only, colback=codebg, colframe=codebg, rounded corners}
int* p = new int[100];
float* f = reinterpret_cast<float*>(p);

int& ri = p[10];
float& rf = reinterpret_cast<float&>(ri);

std::uint64_t address = reinterpret_cast<std::uint64_t>(p);
int* q = reinterpret_cast<int*>(address);
\end{tcblisting}
\pause
\item Обычно нужен в низкоуровневом коде
\pause
\item Нужно использовать с осторожностью
\end{itemize}
\end{frame}

\begin{frame}[fragile]
\frametitle{const\_cast}
\begin{itemize}
\usemintedstyle{tango}
\item Позволяет убрать/добавить константность у указателя или ссылки
\pause
\usemintedstyle{lightbulb}
\begin{tcblisting}{listing engine=minted, minted language=cpp, minted options={fontsize=\scriptsize}, listing only, colback=codebg, colframe=codebg, rounded corners}
const int x = 42;
int* p = &x; // Ошибка компиляции
int* q = const_cast<int*>(&x); // OK
*q = 67; // Undefined behavior
\end{tcblisting}
\pause
\item Нужно использовать с крайней осторожностью
\pause
\usemintedstyle{tango}
\item Почти никогда не нужен, исключение -- реализация парных \mintinline{cpp}|const| и не-\mintinline{cpp}|const| методов у классов
\end{itemize}
\end{frame}

\begin{frame}[fragile]
\frametitle{const\_cast}
\begin{itemize}
\usemintedstyle{tango}
\item Часто, класс предоставляет две версии метода (константную и неконстантную), которые практически не отличаются реализацией
\pause
\usemintedstyle{lightbulb}
\begin{tcblisting}{listing engine=minted, minted language=cpp, minted options={fontsize=\scriptsize}, listing only, colback=codebg, colframe=codebg, rounded corners}
template <typename Key, typename Value>
class map {
  // ...

  const Value& at(const Key& key) const {
    // Поиск в дереве
  }

  Value& at(const Key& key) {
    // Точно такой же поиск в дереве :)
  }
};
\end{tcblisting}
\end{itemize}
\end{frame}

\begin{frame}[fragile]
\frametitle{const\_cast}
\begin{itemize}
\usemintedstyle{tango}
\item Возможное решение -- использовать \mintinline{cpp}|const_cast|
\pause
\usemintedstyle{lightbulb}
\begin{tcblisting}{listing engine=minted, minted language=cpp, minted options={fontsize=\scriptsize}, listing only, colback=codebg, colframe=codebg, rounded corners}
template <typename Key, typename Value>
class map {
  // ...

  const Value& at(const Key& key) const {
    // Поиск в дереве
  }

  Value& at(const Key& key) {
    // Здесь const_cast безопасен: мы знаем, что объект *this
    // на самом деле неконстантный
    return const_cast<Value&>(
      const_cast<const map&>(*this).at(key)
    );
  }
};
\end{tcblisting}
\end{itemize}
\end{frame}

\begin{frame}[fragile]
\frametitle{dynamic\_cast}
\begin{itemize}
\usemintedstyle{tango}
\item Безопасное приведение указателя/ссылки на базовый класс к указателю на родительский класс
\pause
\item Если приведение некорректно, возвращается нулевой указатель, либо (в случае приведения типов ссылок) бросается исключение \mintinline{cpp}|std::bad_cast|
\pause
\usemintedstyle{lightbulb}
\begin{tcblisting}{listing engine=minted, minted language=cpp, minted options={fontsize=\scriptsize}, listing only, colback=codebg, colframe=codebg, rounded corners}
struct A {
  virtual ~A() {}
};

struct B : A { ... };

struct C : A { ... };

A * pa = new B;

dynamic_cast<B*>(pa); // вернёт реальный указатель
dynamic_cast<C*>(pa); // вернёт нулевой указатель

dynamic_cast<B&>(*pa); // вернёт реальную ссылку
dynamic_cast<C&>(*pa); // бросит std::bad_cast
\end{tcblisting}
\end{itemize}
\end{frame}

\begin{frame}[fragile]
\frametitle{Run-Time Type Information (RTTI)}
\begin{itemize}
\usemintedstyle{tango}
\item Как работает \mintinline{cpp}|dynamic_cast|?
\pause
\item Переходит по указателю на таблицу виртуальных функций, рядом с ней компилятор кладёт метаинформацию о типе
\pause
\item \(\Longrightarrow\) \mintinline{cpp}|dynamic_cast| работает только для классов с виртуальными методами
\pause
\item Такая технология называется \textit{Run-Time Type Information (RTTI)}
\pause
\item При программировании под микроконтроллеры или встраиваемые устройства её часто отключают (у GCC флаг \mintinline{text}|-fno-rtti|)
\end{itemize}
\end{frame}

\begin{frame}[fragile]
\frametitle{typeid}
\begin{itemize}
\usemintedstyle{tango}
\item В стандартной библиотеке есть специальный класс \mintinline{cpp}|std::type_info| в файле \mintinline{cpp}|<typeinfo>|
\pause
\item Он позволяет узнать немного информации о типе объекта -- вывести имя \textit{(после name mangling'а!)} или посчитать хеш
\pause
\item Чтобы достать \mintinline{cpp}|std::type_info| по конкретному объекту или типу, есть оператор \mintinline{cpp}|typeid|
\pause
\usemintedstyle{lightbulb}
\begin{tcblisting}{listing engine=minted, minted language=cpp, minted options={fontsize=\scriptsize}, listing only, colback=codebg, colframe=codebg, rounded corners}
#include <typeinfo>
#include <iostream>

int main() {
  int i = 10;
  std::cout << typeid(i).name() << std::endl;
  std::cout << typeid(std::ostream&).name() << std::endl;
}
\end{tcblisting}
\end{itemize}
\end{frame}

\begin{frame}[fragile]
\frametitle{typeid + RTTI}
\begin{itemize}
\usemintedstyle{tango}
\item Если не выключен RTTI, и тип имеет виртуальные функции, то \mintinline{cpp}|typeid| будет использовать RTTI и вернёт настоящий тип объекта:
\usemintedstyle{lightbulb}
\begin{tcblisting}{listing engine=minted, minted language=cpp, minted options={fontsize=\scriptsize}, listing only, colback=codebg, colframe=codebg, rounded corners}
struct A {
  virtual ~A() {}
};

struct B : A {};

A * pa = new B;
typeid(*pa); // вернёт typeid(B)
\end{tcblisting}
\end{itemize}
\end{frame}

\begin{frame}[fragile]
\frametitle{type\_index}
\begin{itemize}
\usemintedstyle{tango}
\item \mintinline{cpp}|std::type_info| -- неудобный (некопируемый) тип с реализацией, зависящей от компилятора
\pause
\item Чтобы было удобнее хранить такие объекты в коллекциях, есть класс-обёртка \mintinline{cpp}|std::type_index|
\pause
\item Фактически, он просто хранит указатель на \mintinline{cpp}|std::type_info|
\pause
\usemintedstyle{lightbulb}
\begin{tcblisting}{listing engine=minted, minted language=cpp, minted options={fontsize=\scriptsize}, listing only, colback=codebg, colframe=codebg, rounded corners}
std::set<std::type_index> set_of_types;
set_of_types.insert(typeid(int));
set_of_types.insert(typeid(void*));
set_of_types.insert(typeid(std::string));
\end{tcblisting}
\end{itemize}
\end{frame}

\begin{frame}[fragile]
\frametitle{Name mangling}
\begin{itemize}
\usemintedstyle{tango}
\item Если хотите получить нормальное имя типа (до name mangling'а, т.н. \textit{demangled name}), в библиотеке Boost есть функция, которая это делает:
\pause
\usemintedstyle{lightbulb}
\begin{tcblisting}{listing engine=minted, minted language=cpp, minted options={fontsize=\scriptsize}, listing only, colback=codebg, colframe=codebg, rounded corners}
#include <boost/core/demangle.hpp>
#include <iostream>

int main() {
  using namespace boost::core;
  std::cout << demangle(typeid(int).name()) << std::endl;
  std::cout << demangle(typeid(std::string).name()) << std::endl;
}
\end{tcblisting}
\end{itemize}
\end{frame}

\begin{frame}[fragile]
\frametitle{lvalue vs rvalue}
\begin{itemize}
\usemintedstyle{tango}
\item Все выражения в C++ делятся на два типа: те, кому можно что-то присвоить, и те, кому нельзя
\pause
\item Те, кому можно присвоить, могут стоять \textit{слева} от оператора присваивания, поэтому называются \textbf{lvalue} (left-value)
\pause
\item Те, кому нельзя, могут стоять только \textit{справа} от оператора присваивания, поэтому называются \textbf{rvalue} (right-value)
\end{itemize}
\end{frame}

\begin{frame}[fragile]
\frametitle{lvalue vs rvalue}
\begin{itemize}
\usemintedstyle{tango}
\item lvalue:
\pause
\begin{itemize}
\item Именованные переменные, разыменовывание указателя, обращение по индексу массива, оращение по ссылке
\pause
\item Существуют и вне выражения, в котором использованы
\pause
\item Имеют адрес в памяти
\end{itemize}
\pause
\item rvalue:
\pause
\begin{itemize}
\item Временные значения
\pause
\item Уничтожаются после вычисления выражения
\end{itemize}
\pause
\item \textbf{NB}: точные правила много сложнее, а ещё есть \textit{xvalue}, \textit{glvalue} и \textit{prvalue}, но их знание не нужно в 99.99\% случаев
\end{itemize}
\end{frame}

\begin{frame}[fragile]
\frametitle{lvalue vs rvalue}
\begin{itemize}
\usemintedstyle{lightbulb}
\begin{tcblisting}{listing engine=minted, minted language=cpp, minted options={fontsize=\scriptsize}, listing only, colback=codebg, colframe=codebg, rounded corners}
int i;

// i - lvalue, 42 - rvalue
i = 42;

int * p = ...;
// p - lvalue, *p - lvalue, i - lvalue
*p = i;

// p[10] - lvalue, p[2] - lvalue, p[2] + 1 - rvalue
p[10] = p[2] + 1;

//    p + 10   rvalue
// *(p + 10)   lvalue
//         i   lvalue
//     i - 5   rvalue
*(p + 10) = i - 5;
\end{tcblisting}
\end{itemize}
\end{frame}

\begin{frame}[fragile]
\frametitle{Move семантика, rvalue-ссылки}
\begin{itemize}
\usemintedstyle{tango}
\item Часто бывают типы, объекты которых не хочется или невозможно копировать: большие динамические массивы, файлы, сетевые сокеты, обёртки других объектов, и т.п.
\pause
\item Даже если копирование можно реализовать, оно влечёт накладные расходы
\pause
\usemintedstyle{lightbulb}
\begin{tcblisting}{listing engine=minted, minted language=cpp, minted options={fontsize=\scriptsize}, listing only, colback=codebg, colframe=codebg, rounded corners}
vector<int> generate();
void consume(vector<int> v);

int main() {
  // 1. Создали временный массив в возвращаемом значении
  // 2. Скопировали временный массив в v
  // 3. Удалили временный массив
  vector<int> v = generate();

  // Скопировали v в аргумент функции
  consume(v);

  // Удалили v
}
\end{tcblisting}
\end{itemize}
\end{frame}

\begin{frame}[fragile]
\frametitle{Move семантика, rvalue-ссылки}
\begin{itemize}
\usemintedstyle{tango}
\item Хочется уметь \textit{передавать владение} объектом из одного места кода в другое
\pause
\item Решение: добавим новый тип ссылок -- \textit{rvalue-ссылки} \mintinline{cpp}|T&&|
\pause
\item Такая ссылка обычно означает временное значение, которое скоро удалится, и у которого можно забрать данные
\pause
\item При передаче rvalue объектов в функции, они предпочитают перегрузки, принимающие rvalue-ссылку
\end{itemize}
\end{frame}

\begin{frame}[fragile]
\frametitle{Move семантика, rvalue-ссылки}
\begin{itemize}
\usemintedstyle{lightbulb}
\begin{tcblisting}{listing engine=minted, minted language=cpp, minted options={fontsize=\scriptsize}, listing only, colback=codebg, colframe=codebg, rounded corners}
class Array {
private:
  int * data;
  size_t size;

public:
  Array(): data(nullptr), size(0) {}
  Array(size_t size) : data(new int[size]), size(size) {}

  // Обычное копирование
  Array(const Array& other)
    : data(new int[other.size]), size(other.size)
  { /* std::copy(...) */ }

  // Конструктор от rvalue-ссылки => можем "забрать себе" данные
  Array(Array&& other)
    : data(other.data), size(other.size)
  {
    // Чтобы избежать double free
    other.data = nullptr; other.size = 0;
  }
};
\end{tcblisting}
\end{itemize}
\end{frame}

\begin{frame}[fragile]
\frametitle{Move семантика, rvalue-ссылки}
\begin{itemize}
\usemintedstyle{lightbulb}
\begin{tcblisting}{listing engine=minted, minted language=cpp, minted options={fontsize=\scriptsize}, listing only, colback=codebg, colframe=codebg, rounded corners}
Array a1(100);

// Вызов обычного копирования
Array a2(a1);

// Вызов конструктора от rvalue объекта
Array a3(A(200));
Array a4(makeArray());
\end{tcblisting}
\end{itemize}
\end{frame}

\begin{frame}[fragile]
\frametitle{Move семантика, rvalue-ссылки}
\begin{itemize}
\usemintedstyle{tango}
\item Обычно вместе с конструктором пишут оператор присваивания от rvalue:
\usemintedstyle{lightbulb}
\begin{tcblisting}{listing engine=minted, minted language=cpp, minted options={fontsize=\scriptsize}, listing only, colback=codebg, colframe=codebg, rounded corners}
class Array {
...

  Array& operator=(Array&& other) {
    if (this == &other) return *this;
    delete data;
    data = other.data;
    size = other.size;
    other.data = nullptr;
    other.size = nullptr;
    return *this;
  }
};
\end{tcblisting}
\pause
\item В общем случае rvalue-ссылки можно использовать где угодно: как типы переменных, как параметры функций, и т.п.
\pause
\item Обычно они имеют смысл только в специализированных конструкторах и операторах присваивания
\end{itemize}
\end{frame}

\begin{frame}[fragile]
\frametitle{Move семантика, rvalue-ссылки}
\begin{itemize}
\usemintedstyle{tango}
\item Что, если мы хотим переложить \textit{(move)} один объект в другой без копирования?
\pause
\item Стандартная функция \mintinline{cpp}|std::move| делает приведение типа из обычной ссылки \mintinline{cpp}|T&| в rvalue-ссылку \mintinline{cpp}|T&&|
\pause
\usemintedstyle{lightbulb}
\begin{tcblisting}{listing engine=minted, minted language=cpp, minted options={fontsize=\scriptsize}, listing only, colback=codebg, colframe=codebg, rounded corners}
Array a1;
Array a2(100);

// Вызовет Array::operator=(Array&&)
a1 = std::move(a2);
\end{tcblisting}
\end{itemize}
\end{frame}

\begin{frame}[fragile]
\frametitle{Move-only объекты}
\begin{itemize}
\usemintedstyle{tango}
\item У класса можно удалить конструктор и оператор обычного копирования, и оставить только rvalue-версии
\pause
\item Такой класс нельзя скопировать, можно только переместить \textit{(move)} значение из одного объекта в другой
\pause
\item Такие классы называют \textit{move-only} или \textit{movable-only}
\pause
\usemintedstyle{lightbulb}
\begin{tcblisting}{listing engine=minted, minted language=cpp, minted options={fontsize=\scriptsize}, listing only, colback=codebg, colframe=codebg, rounded corners}
class Array {
public:
  // Запрет копирования
  Array(const Array&) = delete;
  Array& operator=(const Array& other) = delete;

  // Move-методы
  Array(Array&&);
  Array& operator=(Array&& other);
};
\end{tcblisting}
\end{itemize}
\end{frame}

\begin{frame}[fragile]
\frametitle{Move семантика, rvalue-ссылки}
\begin{itemize}
\usemintedstyle{tango}
\item Возвращаемое значение функции -- автоматически rvalue, поэтому в \mintinline{cpp}|return ...;| не нужно делать \mintinline{cpp}|std::move| (это ещё и испортит RVO/NRVO)
\usemintedstyle{lightbulb}
\begin{tcblisting}{listing engine=minted, minted language=cpp, minted options={fontsize=\scriptsize}, listing only, colback=codebg, colframe=codebg, rounded corners}
Array makeArray() {
  Array a(100);
  // ...

  // Автоматически вызовет rvalue-конструктор, если
  // не сработало NRVO
  return a; 
}
\end{tcblisting}
\end{itemize}
\end{frame}

\begin{frame}[fragile]
\frametitle{Move семантика, rvalue-ссылки}
\begin{itemize}
\usemintedstyle{tango}
\item \mintinline{cpp}|std::swap| на самом деле реализован через \mintinline{cpp}|std::move|, так что его безопасно использовать и для тяжело копируемых объектов (напр. \mintinline{cpp}|std::vector|):
\usemintedstyle{lightbulb}
\begin{tcblisting}{listing engine=minted, minted language=cpp, minted options={fontsize=\scriptsize}, listing only, colback=codebg, colframe=codebg, rounded corners}
template <typename T>
void swap(T& a, T& b) {
  T temp = std::move(a);
  a = std::move(b);
  b = std::move(temp);
}
\end{tcblisting}
\end{itemize}
\end{frame}

\begin{frame}[fragile]
\frametitle{Move семантика vs линейные/аффинные типы}
\begin{itemize}
\usemintedstyle{tango}
\item Весь механизм с использованием rvalue-ссылок называется \textit{move-семантикой}
\pause
\item Линейный/аффинный тип -- термин из теории типов, означающий объекты, которые нужно использовать ровно один раз/не более одного раза
\pause
\item Относительно новая идея для языков программрования (один из примеров -- Rust)
\pause
\item Должны поддерживаться на уровне компилятора
\pause
\item C++ не поддерживает в явном виде линейные или аффинные типы, но позволяет \textit{эмулировать} аффинные типы с помощью move-семантики
\end{itemize}
\end{frame}

\begin{frame}[fragile]
\frametitle{Universal references, perfect forwarding}
\begin{itemize}
\usemintedstyle{tango}
\item Есть случай, когда \mintinline{cpp}|T&&| не означает rvalue-ссылку: если \mintinline{cpp}|T| -- шаблонный параметр функции
\pause
\item Тогда есть особые правила \textit{(reference collapsing rules)}, выводящие \mintinline{cpp}|T| как обычную или rvalue-ссылку
\pause
\item В этом случае \mintinline{cpp}|T&&| называется \textit{универсальной ссылкой (universal reference)}
\pause
\item Вместе со стандартной функцией \mintinline{cpp}|std::forward| позволяет принять неизвестный набор параметров и передать в другую функцию, не потеряв информацию об lvalue/rvalue
\pause
\item Полезно для методов-фабрик, callback'ов, обёрток, и т.п.
\pause
\item Сам механизм называется \textit{perfect forwarding}
\end{itemize}
\end{frame}

\begin{frame}[fragile]
\frametitle{Universal references, perfect forwarding}
\begin{itemize}
\usemintedstyle{lightbulb}
\begin{tcblisting}{listing engine=minted, minted language=cpp, minted options={fontsize=\scriptsize}, listing only, colback=codebg, colframe=codebg, rounded corners}
template <typename T>
void foo(T&& x) {
  bar(std::forward<T>(x));
}
\end{tcblisting}
\pause
\item Часто используется вместе с variadic templates:
\usemintedstyle{lightbulb}
\begin{tcblisting}{listing engine=minted, minted language=cpp, minted options={fontsize=\scriptsize}, listing only, colback=codebg, colframe=codebg, rounded corners}
template <typename T, typename ... Args>
std::shared_ptr<T> std::make_shared(Args&& ... args) {
  return std::shared_ptr<T>(
    new T(std::forward<Args>(args)...);
  );
}
\end{tcblisting}
\end{itemize}
\end{frame}

\begin{frame}[fragile]
\frametitle{Вывод типов}
\begin{itemize}
\usemintedstyle{tango}
\item Часто не хочется писать слишком длинный тип переменной, или его тяжело получить
\pause
\usemintedstyle{lightbulb}
\begin{tcblisting}{listing engine=minted, minted language=cpp, minted options={fontsize=\scriptsize}, listing only, colback=codebg, colframe=codebg, rounded corners}
for (std::map<std::string, std::vector<int>>::const_iterator
  it = m.begin(); it != m.end(); ++it) { ... }
\end{tcblisting}
\pause
\usemintedstyle{lightbulb}
\begin{tcblisting}{listing engine=minted, minted language=cpp, minted options={fontsize=\scriptsize}, listing only, colback=codebg, colframe=codebg, rounded corners}
template <typename Iterator>
void qsort(Iterator being, Iterator end) {
  typename std::iterator_traits<Iterator>::value_type
    pivot = ...;
  ...
}
\end{tcblisting}
\pause
\usemintedstyle{tango}
\item Решение: \mintinline{cpp}|auto|
\end{itemize}
\end{frame}

\begin{frame}[fragile]
\frametitle{Вывод типов}
\begin{itemize}
\usemintedstyle{tango}
\item Ключевое слово \mintinline{cpp}|auto| можно использовать вместо типа переменной, тогда тип выведется автоматически из выражения, которым инициализируется переменная
\pause
\usemintedstyle{lightbulb}
\begin{tcblisting}{listing engine=minted, minted language=cpp, minted options={fontsize=\scriptsize}, listing only, colback=codebg, colframe=codebg, rounded corners}
// int
auto i = 10;

// const char*
auto str = "Hello, world!";

std::map<std::string, std::vector<int>> m;

// std::map<std::string, std::vector<int>>::iterator
auto it = m.begin();

// std::pair<const std::string, std::vector<int>>
// NB: не ссылка, а копия!
auto pivot = *it;
\end{tcblisting}
\end{itemize}
\end{frame}

\begin{frame}[fragile]
\frametitle{Вывод типов}
\begin{itemize}
\usemintedstyle{tango}
\item К \mintinline{cpp}|auto| можно добавить спецификаторы, уточняющие тип:
\pause
\usemintedstyle{lightbulb}
\begin{tcblisting}{listing engine=minted, minted language=cpp, minted options={fontsize=\scriptsize}, listing only, colback=codebg, colframe=codebg, rounded corners}
int * p = ...;

// int
auto x = *p;

// const int
const auto x = *p;

// int&
auto& x = *p;
\end{tcblisting}
\end{itemize}
\end{frame}

\begin{frame}[fragile]
\frametitle{Вывод типов}
\begin{itemize}
\usemintedstyle{tango}
\item Оператор \mintinline{cpp}|decltype(x)| подставляет вместо себя тип выражения \mintinline{cpp}|(x)|
\pause
\item Можно использовать везде, где ожидается тип
\pause
\usemintedstyle{lightbulb}
\begin{tcblisting}{listing engine=minted, minted language=cpp, minted options={fontsize=\scriptsize}, listing only, colback=codebg, colframe=codebg, rounded corners}
// int
decltype(10) i = 20;

// float
decltype(1 + 2.f) f = 3.14f;

int* p = ...;

// int*
decltype(p + 10) q = ...;

// int&
decltype(p[10]) r = ...;
\end{tcblisting}
\pause
\usemintedstyle{tango}
\item \textbf{NB}: есть фокусы связанные со ссылками, см. \mintinline{cpp}|decltype(auto)| и \mintinline{cpp}|decltype((x))|
\end{itemize}
\end{frame}

\begin{frame}[fragile]
\frametitle{Вывод типов}
\begin{itemize}
\usemintedstyle{tango}
\item \mintinline{cpp}|decltype| часто полезен в сложном шаблонном коде
\pause
\item \mintinline{cpp}|auto| нужно использовать аккуратно
\pause
\item Где-то \mintinline{cpp}|auto| упрощает чтение кода (напр. с итераторами -- мы и так знаем, что там за тип)
\pause
\item Где-то \mintinline{cpp}|auto| усложняет чтение кода (непонятно, какая переменная в итоге какого типа)
\end{itemize}
\end{frame}

\begin{frame}[fragile]
\frametitle{Range-based for loop}
\begin{itemize}
\usemintedstyle{tango}
\item Даже с итераторами, простая концепция `пройтись по всем элементам коллекции' выглядит довольно громоздко:
\usemintedstyle{lightbulb}
\begin{tcblisting}{listing engine=minted, minted language=cpp, minted options={fontsize=\scriptsize}, listing only, colback=codebg, colframe=codebg, rounded corners}
for (auto it = v.begin(); it != v.end(); ++it) {
  int x = *it;
  // ...
}
\end{tcblisting}
\pause
\item Решение: \textit{range-based for loop}:
\usemintedstyle{lightbulb}
\begin{tcblisting}{listing engine=minted, minted language=cpp, minted options={fontsize=\scriptsize}, listing only, colback=codebg, colframe=codebg, rounded corners}
for (int x : v) {
  // ...
}
\end{tcblisting}
\pause
\item Превращается компилятором в точно такой же код с итераторами
\end{itemize}
\end{frame}

\begin{frame}[fragile]
\frametitle{Range-based for loop}
\begin{itemize}
\usemintedstyle{tango}
\item Как при объявлении переменной, вместо \mintinline{cpp}|int| можно написать \mintinline{cpp}|auto| или принимать по ссылке:
\usemintedstyle{lightbulb}
\begin{tcblisting}{listing engine=minted, minted language=cpp, minted options={fontsize=\scriptsize}, listing only, colback=codebg, colframe=codebg, rounded corners}
std::vector<int> v;

// Локальный x - копия элемента коллекции
for (int x : v) { ... }

// Тип выведется как int
for (auto x : v) { ... }

// Локальный x - ссылка на элемент коллекции
for (int& x : v) { ... }

// Тип выведется как int&
for (auto& x : v) { ... }
\end{tcblisting}
\end{itemize}
\end{frame}

\begin{frame}[fragile]
\frametitle{Range-based for loop}
\begin{itemize}
\usemintedstyle{tango}
\item Если бы такой цикл использовал \textit{методы} \mintinline{cpp}|begin()| и \mintinline{cpp}|end()|, то он бы не работал с обычными массивами:
\usemintedstyle{lightbulb}
\begin{tcblisting}{listing engine=minted, minted language=cpp, minted options={fontsize=\scriptsize}, listing only, colback=codebg, colframe=codebg, rounded corners}
int array[5];
for (int v : array) { ... }
\end{tcblisting}
\pause
\usemintedstyle{tango}
\item На самом деле, вместо этого range-based for использует стандартные функции \mintinline{cpp}|std::begin()| И \mintinline{cpp}|std::end()|
\pause
\item Для обычных массивов они возвращают нужные указатели, а иначе вызывают метод \mintinline{cpp}|v.begin()| или \mintinline{cpp}|v.end()|
\end{itemize}
\end{frame}

\end{document}